\documentclass{article}

\usepackage{tabularx}
\usepackage{booktabs}
\usepackage{float}
\usepackage{longtable}


\title{Reflection and Traceability Report on \progname}

\author{\authname}

\date{}

\input{../Comments}
\input{../Common}

\begin{document}

\maketitle

\plt{Reflection is an important component of getting the full benefits from a
learning experience.  Besides the intrinsic benefits of reflection, this
document will be used to help the TAs grade how well your team responded to
feedback.  Therefore, traceability between Revision 0 and Revision 1 is and
important part of the reflection exercise.  In addition, several CEAB (Canadian
Engineering Accreditation Board) Learning Outcomes (LOs) will be assessed based
on your reflections.}

\section{Changes in Response to Feedback}

\plt{Summarize the changes made over the course of the project in response to
feedback from TAs, the instructor, teammates, other teams, the project
supervisor (if present), and from user testers.}

\plt{For those teams with an external supervisor, please highlight how the feedback 
from the supervisor shaped your project.  In particular, you should highlight the 
supervisor's response to your Rev 0 demonstration to them.}

\plt{Version control can make the summary relatively easy, if you used issues
and meaningful commits.  If you feedback is in an issue, and you responded in
the issue tracker, you can point to the issue as part of explaining your
changes.  If addressing the issue required changes to code or documentation, you
can point to the specific commit that made the changes.  Although the links are
helpful for the details, you should include a label for each item of feedback so
that the reader has an idea of what each item is about without the need to click
on everything to find out.}

\plt{If you were not organized with your commits, traceability between feedback
and commits will not be feasible to capture after the fact.  You will instead
need to spend time writing down a summary of the changes made in response to
each item of feedback.}

\plt{You should address EVERY item of feedback.  A table or itemized list is
recommended.  You should record every item of feedback, along with the source of
that feedback and the change you made in response to that feedback.  The
response can be a change to your documentation, code, or development process.
The response can also be the reason why no changes were made in response to the
feedback.  To make this information manageable, you will record the feedback and
response separately for each deliverable in the sections that follow.}

\plt{If the feedback is general or incomplete, the TA (or instructor) will not
be able to grade your response to feedback.  In that case your grade on this
document, and likely the Revision 1 versions of the other documents will be
low.} 

\subsection{SRS and Hazard Analysis}
\subsubsection{Hazard Analysis}

\begin{longtable}{|>{\raggedright\arraybackslash}p{0.24\textwidth}|>{\raggedright\arraybackslash}p{0.30\textwidth}|>{\raggedright\arraybackslash}p{0.40\textwidth}|}
\caption{Hazard Analysis Feedback Traceability}
\label{tab:hazard-feedback-traceability}\\
\hline
\textbf{Feedback Source} & \textbf{Feedback Summary} & \textbf{Changes Made (or justification why not)} \\
\hline
\endfirsthead

\hline
\textbf{Feedback Source} & \textbf{Feedback Summary} & \textbf{Changes Made (or justification why not)} \\
\hline
\endhead

\hline
\endfoot

\hline
\endlastfoot

Peer Review (\href{https://github.com/VirochaanRG/MES-Event-Management-System/issues/67}{Issue \#67}) & Provide Metrics for Hazards for better clarity of the severity and nature & Added a column in our FMEA table which includes the severity which we colour coded by importance. This was a worthwhile change as it adds better clarity to the user on the level of threat\\
\hline
Peer Review (\href{https://github.com/VirochaanRG/MES-Event-Management-System/issues/73}{Issue \#73}) &  Update Scope and purpose to be more specific and detailed & Section 2.2 was made more defined and detailed. Short description of the losses that may have been incurred is also included within.\\
\hline
Peer Review (\href{https://github.com/VirochaanRG/MES-Event-Management-System/issues/68}{Issue \#68}) & Turn the critical assumption of malicious usage into hazard & \textbf{Not Implemented:} We considered this a reasonable assumption as the goal of the hazard analysis was not to stop all malicious actors from messing with the system but rather ensure that the system functioned correctly and did not provide an easy entry or failure mode in which it could be exploited. For the current scope of our project it is unreasonable to be able to defend on all potential malicious attacks by users and is something that we can look towards working on in the future.\\
\hline
TA Feedback (\href{https://github.com/VirochaanRG/MES-Event-Management-System/issues/309}{Issue \#309}) & Add a system boudary diagram to better understand of the components interact  & Added section 3.3 with a system diagram showcasing the interactions between the various components within our system in relevance the the hazard analysis. Provides users with a better understanding on the fundamentals of how this system works.\\
\hline
\end{longtable}

\subsection{Design and Design Documentation}

\subsection{VnV Plan and Report}
\subsubsection{VnV Plan}

\begin{longtable}{|>{\raggedright\arraybackslash}p{0.24\textwidth}|>{\raggedright\arraybackslash}p{0.30\textwidth}|>{\raggedright\arraybackslash}p{0.40\textwidth}|}
\caption{VnV Plan Feedback Traceability}
\label{tab:vnv-feedback-traceability}\\
\hline
\textbf{Feedback Source} & \textbf{Feedback Summary} & \textbf{Changes Made (or justification why not)} \\
\hline
\endfirsthead

\hline
\textbf{Feedback Source} & \textbf{Feedback Summary} & \textbf{Changes Made (or justification why not)} \\
\hline
\endhead

\hline
\endfoot

\hline
\endlastfoot
Peer Review (\href{https://github.com/VirochaanRG/MES-Event-Management-System/issues/84}{Issue \#84}) & Fix Spelling and Grammar Issues & Ran spellcheck throughout the document and just ensured everything made sense.\\
\hline
Peer Review (\href{https://github.com/VirochaanRG/MES-Event-Management-System/issues/87}{Issue \#87}) & Make system Tests less ambiguous & In section 4.1, every single functional system test case was looked at and improved to be be a bit more specific with the testing methodology and thoroughness of the testing. We also looked at making the output more verifiable by looking directly at database rather than visual confirmations.\\
\hline
Peer Review (\href{https://github.com/VirochaanRG/MES-Event-Management-System/issues/85}{Issue \#85}) & Too Detailed Verification of SRS& \textbf{Not Implemented:} We considered that the requirements validation as part of the verification of the SRS as included within the documentation and the rubric of said document.\\
\hline
Peer Review (\href{https://github.com/VirochaanRG/MES-Event-Management-System/issues/88}{Issue \#88}) & Untested Non-functional requirements& In section 4.2, we went over the system tests to ensure that all the requirements in our SRS were being covered by the current tests. The only one missing was SC-1 so we added a new test covering that. We also reworked existing tests to ensure that they were better covering the NFRs as required.\\
\hline
Peer Review (\href{https://github.com/VirochaanRG/MES-Event-Management-System/issues/93}{Issue \#93}) & Unclear Criterion for PF-2 & In section 4.2.4 we included a specific Measured time to be less than 4 seconds. We also looked at the rest of tests and ensured proper criterions were established as neccesary\\
\hline
Peer Review (\href{https://github.com/VirochaanRG/MES-Event-Management-System/issues/92}{Issue \#92}) & Section 4.3 Table could be more detailed & \textbf{Not Implemented:} We didn't feel it neccesary to include all the details again and felt it would make the table more cluttered and didn't actually fit the nature of a traceability table. As a compromise we once again included the link to our SRS document were the readers could gain more context if necessary.\\
\hline
Peer Review (\href{https://github.com/VirochaanRG/MES-Event-Management-System/issues/89}{Issue \#89}) & Include Usability Survey Questions & We included a link to our usability report extra which has all the survey questions as well as additional context and findings which readers may find useful. \\
\hline
Peer Review + TA Feedback (\href{https://github.com/VirochaanRG/MES-Event-Management-System/issues/86}{Issue \#86})  & NFR Traceability Table & We included a table containing traceability between the NFR and the System Tests to ensure proper connection between the SRS and VNV Plan. This can be seen in table 5 in section 4.3\\
\hline
Peer Review (\href{https://github.com/VirochaanRG/MES-Event-Management-System/issues/95}{Issue \#95})  &SUpervisor Involvement in Verification & \textbf{Not Implemented:} Section 3.4 is meant to be a verification and validation of our internal strategies so we didn't believe supervisor validation was necessary in that situation. Section 3.5 has supervisor verification done through the usability verification as they will be one of the people doing the testing. \\
\hline
\end{longtable}

\section{Extras}

\plt{Summarize the extras that were tackled by this project.  Extras
can include usability testing, code walkthroughs, user documentation, formal
proof, GenderMag personas, Design Thinking, etc.  Extras should have already
been approved by the course instructor as included in your problem statement.}

\section{Design Iteration (LO11 (PrototypeIterate))}

\plt{Explain how you arrived at your final design and implementation.  How did
the design evolve from the first version to the final version?} 

\plt{Don't just say what you changed, say why you changed it.  The needs of the
client should be part of the explanation.  For example, if you made changes in
response to usability testing, explain what the testing found and what changes
it led to.}

\section{Design Decisions (LO12)}

\plt{Reflect and justify your design decisions.  How did limitations,
 assumptions, and constraints influence your decisions?  Discuss each of these
 separately.}

\section{Economic Considerations (LO23)}

\plt{Is there a market for your product? What would be involved in marketing your 
product? What is your estimate of the cost to produce a version that you could 
sell?  What would you charge for your product?  How many units would you have to 
sell to make money? If your product isn't something that would be sold, like an 
open source project, how would you go about attracting users?  How many potential 
users currently exist?}

\section{Reflection on Project Management (LO24)}

\plt{This question focuses on processes and tools used for project management.}

\subsection{How Does Your Project Management Compare to Your Development Plan}

\plt{Did you follow your Development plan, with respect to the team meeting plan, 
team communication plan, team member roles and workflow plan.  Did you use the 
technology you planned on using?}

\subsection{What Went Well?}

\plt{What went well for your project management in terms of processes and 
technology?}

\subsection{What Went Wrong?}

\plt{What went wrong in terms of processes and technology?}

\subsection{What Would you Do Differently Next Time?}

\plt{What will you do differently for your next project?}

\section{Ethics and Equity (LO18)}

\plt{Describe how your team applied the principle of equity and universal design
to ensure equitable treatment of all stakeholders.}

\section{Reflection on Capstone}

\plt{This question focuses on what you learned during the course of the capstone project.}

\subsection{Which Courses Were Relevant}

\plt{Which of the courses you have taken were relevant for the capstone project?}

\subsection{Knowledge/Skills Outside of Courses}

\plt{What skills/knowledge did you need to acquire for your capstone project
that was outside of the courses you took?}

\end{document}